\begin{pastexam}[label={px:ntust-109-ice-06}]{臺科大資工所 109 學年度 第 6 題}
（問答）

An array contains a list of $n$ items whose keys are distinct and happen to be in ascending order
already. Unaware of this fact, we want to sort the list into ascending order. For each of the
following sorting methods, give the corresponding time complexity (in terms of $n$ using $O(\ )$
notation) of the sorting operation.
\begin{enumerate}[label=(\alph*)]
\item Insertion Sort (3\%)
\item Quick Sort (3\%)
\end{enumerate}
\end{pastexam}

\begin{pastsoln}
答案：(a) $O(n)$；(b) $O(n^2)$。

題目沒有說快速排序怎麼選基準，下面照課本最常見的寫法：固定取第一個或最後一個元素當基準
（資結書虛擬碼~14.9 的 Lomuto 分割取 $A[r]$）。 改用隨機選基準或三數取中位數，(b) 就不是這個答案，
最後再說明。

\textbf{(a) 插入排序}。 第 $j$ 輪要把 $A[j]$ 插進已排好的 $A[1..j-1]$（資結書虛擬碼~1.3）。
輸入已經由小到大時，$A[j-1] < A[j]$，內層迴圈比較一次就停，不挪動任何元素。 每一輪 $O(1)$，
共 $n - 1$ 輪，比較 $n - 1$ 次，總時間 $O(n)$（資結書第~1.8~節的最好情況）。

\textbf{(b) 快速排序}。 以 Lomuto 分割、基準取最後一個為例：子陣列 $A[p..r]$ 已經排好，基準
$A[r]$ 是最大值，前面 $r - p$ 個都比它小，分割完基準留在原位，切成 $r - p$ 個和 $0$ 個。 左邊
那段仍然排好，下一層又是同樣的情形，每一層只少一個元素。 第 $k$ 層分割的長度是 $n - k + 1$，
比較 $n - k$ 次，總比較次數
\[
(n - 1) + (n - 2) + \cdots + 1 = \frac{n(n-1)}{2},
\]
遞迴式 $T(n) = T(n-1) + \Theta(n)$，$T(n) = \Theta(n^2)$，是 $O(n^2)$（資結書第~14.5~節「最壞和最好
的情況」）。 基準取第一個元素時也一樣：基準是最小值，每次切成 $0$ 個和 $m - 1$ 個。

\textbf{如果基準選法不同}。 隨機選基準時，期望時間是 $O(n \lg n)$，和輸入順序無關；三數取中位數
（第一個、中間、最後一個）在已排好的輸入上會選到正中間，每次切成兩半，也是 $O(n \lg n)$。 作答時
寫出採用哪一種基準，閱卷才看得出答案從哪裡來。
\end{pastsoln}
